\chapter{Potenzen und Idempotente}

\section{Existenz und Schranken von Index und Periode}

\begin{definition}[Index und Periode]
Für $a\in S$ heißen die Zahlen
\begin{align*}
 \mu(a)&=\min\{r\geq1:\exists d\geq1\ (a^{r+d}=a^r)\},\\
 \lambda(a)&=\min\{d\geq1:a^{\mu(a)+d}=a^{\mu(a)}\}
\end{align*}
\emph{Index} und \emph{Periode} von $a$. Der folgende Satz sichert die
Existenz dieser Minima. Die Anlauflänge ist $\mu(a)-1$.
\end{definition}

\begin{theorem}[Existenz und Schranke von Index und Periode]\label{guide:period}
Für jedes $a\in S$ existieren Index $\mu(a)$ und Periode $\lambda(a)$.
Sie erfüllen
\[
 \mu(a)+\lambda(a)-1\leq n.
\]
\cite{projectindex,harju}
\end{theorem}

\section{Gemeinsame Exponenten und zyklische Untergruppen}

\begin{definition}[Gemeinsamer Index und gemeinsame Periode]
Für $S\ne\varnothing$ setzen wir
\[
 M(S)=\max_{a\in S}\mu(a),\qquad
 P(S)=\operatorname{kgV}_{a\in S}\lambda(a).
\]
Für $S=\varnothing$ setzen wir $M(S)=P(S)=1$.
\end{definition}

\begin{theorem}[Gemeinsame Potenzparameter und Idempotenzexponent]\label{guide:uniform}
Für jede endliche Halbgruppe $S$ existieren die gemeinsamen Parameter
$M(S)$ und $P(S)$, und es gilt
\[
 M(S)\leq\max(1,n),\qquad P(S)\mid n!.
\]
Der Exponent $n!$ macht gleichzeitig alle Potenzen $a^{n!}$ idempotent:
\[
 a^{2n!}=a^{n!}\qquad(a\in S),
\]
wobei $0!=1$ gesetzt wird. \cite{projectindex,pin}
\end{theorem}

\begin{definition}[Gemeinsamer Idempotenzexponent und Omega-Potenz]
Der kleinste gemeinsame Idempotenzexponent ist
\[
 \kappa(S)=\min\{q\geq1:\forall a\in S\ (a^{2q}=a^q)\}.
\]
Für $a\in S$ bezeichnet $a^\omega$ die Potenz $a^{\kappa(S)}$.
Für $j\geq1$ schreiben wir $a^{\omega+j}=a^{\kappa(S)+j}$.
Das Zeichen $\omega$ bezeichnet hier keinen festen natürlichen Exponenten
für alle endlichen Halbgruppen zugleich.
\end{definition}

\begin{theorem}[Zyklische Untergruppen hoher Potenzen]\label{guide:kernelpower}
Für jedes $a\in S$ existiert eine zyklische Untergruppe
$C_a\subseteq\langle a\rangle$ mit
\[
 a^r\in C_a\qquad\text{für alle }r\geq\max(1,n).
\]
Insbesondere gilt $S\ne\varnothing\Rightarrow E(S)\ne\varnothing$.
\cite{projectindex,grillet}
\end{theorem}

\begin{theorem}[Idempotentenfaktorisierung langer Produkte]\label{guide:longproducts}
Für jedes $k\geq\max(1,n)$ gilt
\[
 S^k\subseteq SE(S)S.
\]
\cite[Lemma~2.2]{AlmeidaMargolisSteinbergVolkov2009}
\end{theorem}
